\documentclass[10pt,a4paper]{article}
% Examples are included from the checked Rust sources so signatures stay aligned.
\usepackage[T1]{fontenc}
\usepackage{lmodern,microtype}
\usepackage[margin=22mm]{geometry}
\usepackage{booktabs,tabularx,array,xcolor,listings,enumitem,fancyhdr}
\usepackage[pdftitle={Digital RF in Rust: API guide},pdfauthor={},colorlinks=true,linkcolor=blue!45!black,urlcolor=blue!45!black]{hyperref}
\definecolor{ink}{HTML}{17394D}
\definecolor{wash}{HTML}{F4F6F8}
\lstdefinelanguage{Rust}{keywords={use,fn,let,mut,pub,struct,enum,impl,if,else,for,in,match,return,as,Some,None,Ok,Err,true,false},sensitive=true,morecomment=[l]{//},morestring=[b]"}
\lstset{language=Rust,basicstyle=\ttfamily\footnotesize,keywordstyle=\color{ink}\bfseries,commentstyle=\color{black!55},stringstyle=\color{blue!45!black},backgroundcolor=\color{wash},frame=single,rulecolor=\color{black!12},framesep=6pt,breaklines=true,columns=fullflexible,keepspaces=true,showstringspaces=false,aboveskip=8pt,belowskip=8pt}
\setlength{\parindent}{0pt}\setlength{\parskip}{6pt}
\setlist[itemize]{nosep,leftmargin=1.5em}
\pagestyle{fancy}\fancyhf{}\fancyhead[L]{\small\textcolor{ink}{rdigital\_rf}}\fancyhead[R]{\small API guide}\fancyfoot[C]{\thepage}
\setlength{\headheight}{14pt}
\newcommand{\api}[1]{\texttt{\detokenize{#1}}}
\newcommand{\chapterpage}[1]{\clearpage\section{#1}}
\begin{document}
\begin{center}
{\LARGE\bfseries\textcolor{ink}{Digital RF in Rust}}\\[5pt]
{\large rdigital\_rf: API guide and examples}\\[6pt]
Reader API and optional raw binary container
\end{center}
\section{Overview and installation}
\api{rdigital_rf} is a Rust implementation for accessing time-indexed radio
samples. It reads original Digital RF 2.x HDF5 recordings and an optional raw
binary container through the same reader API. Samples keep their absolute
indices and gaps remain explicit. Recording directories do not have to be
combined or converted before reading.

HDF5 remains the default format. Raw writing is opt-in through
\api{raw::RawWriter}. The raw extension separates the high-rate sample bytes
from HDF5 metadata; the upstream Python DigitalRF reader does not support it.
This independent crate implements a subset of DigitalRF functionality, not the
complete Python API. An ordinary Digital RF HDF5 writer is not implemented.

Install Rust, a C toolchain and the native HDF5 development library. On
Debian/Ubuntu, install \api{libhdf5-dev} and \api{pkg-config}; with Homebrew,
install \api{hdf5} and \api{pkg-config}. The crate is not yet on crates.io.
Add the crate from Git:
\begin{lstlisting}[language={}]
[dependencies]
rdigital_rf = { git = "https://github.com/jvierine/rdigital_rf" }
anyhow = "1"
\end{lstlisting}
\subsection*{A first read}
\begin{lstlisting}
use rdigital_rf::Reader;

fn main() -> anyhow::Result<()> {
    let mut r = Reader::open(["/archive/run-a", "/archive/run-b"])?;
    if let Some((first, last)) = r.get_bounds("ch0")? {
        let count = (last - first).min(1023) as usize + 1;
        let iq = r.read_complex32(first, count, "ch0", 0)?;
        println!("{} samples", iq.len());
    }
    Ok(())
}
\end{lstlisting}
\api{Reader::open} accepts a recording root containing channel directories, a
root containing \api{rf_data/}, or an individual channel directory. It discovers
channel layouts at opening; it is not a directory-watching service. Channel
names are directory names, such as \api{ch0}.

\chapterpage{Bounded streaming reads}
Use \api{read_complex32} for a contiguous vector of complex 32-bit floating-point
samples. Its arguments are the absolute first sample, number of samples,
channel name and zero-based subchannel. Integer sample values are widened
without normalization; digital or analog gain calibration belongs to the caller.
\lstinputlisting{../examples/stream_read.rs}
The example accepts a recording root and optional channel name as command-line
arguments. It requests at most 1024 samples and handles missing data separately
from other errors. A later valid read remains possible after a gap error.

\textbf{Bounds are not a continuity guarantee.} Both \api{bounds} and
\api{get_bounds} return \api{Option<(u64, u64)>}: inclusive first and last stored
indices, or \api{None} if no finalized samples are found. Missing files or
indexed gaps can occur between those indices. The strict methods report
\api{MissingSamples}, whose \api{gaps} field contains half-open ranges.

\textbf{Cache limits are not output limits.} The example limits the number of
open files and the HDF5 raw chunk cache per file. The requested output vector,
conversion buffers, file metadata and OS page cache are additional memory.
Bound each read length in a streaming application. Native complex i8/i16 HDF5
samples use direct conversion to complex-f32; raw payloads bypass HDF5 sample
conversion altogether.

\chapterpage{Gap-aware reads and the sample model}
Use \api{read} when processing available blocks is preferable to rejecting the
whole request. It takes a half-open interval: \api{start..end} includes
\api{start}, excludes \api{end}, and contains \api{end - start} samples.
\lstinputlisting{../examples/gap_read.rs}
The example accepts a root, start and end. Indices are unsigned integers, not
floating-point timestamps. If the rational sample rate is $N/D$ samples per
second, an index $s$ corresponds to $sD/N$ seconds after the Unix epoch.
Keep absolute indices as \api{u64}; large indices cannot always be represented
exactly as \api{f64}.

\par\medskip
\begin{tabularx}{\linewidth}{@{}lX@{}}\toprule
Result & Meaning\\\midrule
\api{Read.requested} & Requested half-open sample interval.\\
\api{Read.blocks} & Ordered \api{SampleBlock}s, each with an absolute \api{start}
and a \api{Vec<Complex64>} of available samples.\\
\api{Read.gaps} & Exact half-open missing intervals; never filled with zeros.\\
\api{Read.sources} & \api{SourceSpan}s with sample interval and source metadata-file path.\\
\api{is_complete()} & True when no gaps are present.\\
\api{into_contiguous()} & Concatenated samples for a complete request; otherwise
an error containing \api{MissingSamples}.\\\bottomrule
\end{tabularx}

\api{read_vector} returns \api{Vector { samples, shape }}. With \api{Some(i)},
it selects one subchannel; \api{None} selects all in row-major order
(sample, subchannel). One selected subchannel has shape \api{[length]}; multiple
have shape \api{[length, num_subchannels]}. A one-sample, multiple-subchannel
selection is rejected to match the tested upstream squeeze behavior.

\api{Samples} contains \api{Real32}, \api{Real64}, \api{Complex32} or
\api{Complex64} vectors. Small integer components (up to 16 bits) and float32
promote to float32 or complex-f32; larger supported components promote to
float64 or complex-f64. \api{Vector::into_complex64} provides a uniform complex
representation, with zero imaginary parts for real recordings.

\chapterpage{Reader reference}
All fallible operations return \api{anyhow::Result<T>}. A malformed file,
incompatible layout or conflicting overlap is an error, not an empty read.
\par\medskip
\begin{tabularx}{\linewidth}{@{}p{.47\linewidth}X@{}}\toprule
Method & Result and behavior\\\midrule
\api{Reader::open(roots)} & Discover channels across local roots; deduplicate
identical root arguments. Automatically detect the container per chunk.\\[3pt]
\api{channels()} & \api{Vec<&str>} of discovered channel names.\\[3pt]
\api{properties(channel)} & Borrowed \api{Properties} for the channel.\\[3pt]
\api{bounds(channel)}\newline
\api{get_bounds(channel)} & Inclusive first/last stored indices; aliases. Inspect
file indices rather than sample payloads.\\[3pt]
\api{read(channel, range, subchannel)} & Gap-aware \api{Read}; half-open interval,
including empty intervals.\\[3pt]
\api{read_vector(start, length, channel, subchannel)} & Strict promoted
\api{Vector}; optional subchannel. Positive length required.\\[3pt]
\api{read_complex32(start, length, channel, subchannel)} & Strict
\api{Vec<Complex32>}; one subchannel. Positive length required.\\[3pt]
\api{set_cache_options(options)} & Set \api{CacheOptions} and clear open-file cache.\\[3pt]
\api{set_cache_byte_limit(bytes)} & Positive ceiling on HDF5 raw chunk cache per
file; clears open-file cache.\\[3pt]
\api{cache_stats()} & \api{CacheStats}: opens, cache hits and last HDF5 chunk-cache size.\\[3pt]
\api{clear_cache()} & Close cached file handles; keep discovered channels.\\\bottomrule
\end{tabularx}

\subsection*{Channel properties}
\api{Properties} exposes \api{sample_rate_numerator},
\api{sample_rate_denominator}, \api{num_subchannels}, \api{is_complex},
\api{epoch} and \api{format}. \api{Format} describes the promoted output
representation, not the original storage dtype. Rates are reduced to their
rational form. The supported epoch is \api{1970-01-01T00:00:00Z}.

\subsection*{Cache policy}
\api{CacheOptions::default()} uses \api{minimum_bytes = 64 * 1024 * 1024} and
\api{max_open_files = 8}. Both must be positive. Without a byte ceiling, an
HDF5 cache is sized to at least the minimum and the full logical dataset size.
This favors repeated sparse reads of compressed files but is not a total RAM
budget. \api{CacheStats} contains \api{file_opens}, \api{file_hits} and
\api{last_chunk_cache_bytes}; the last field describes HDF5 caching, not raw I/O.

\subsection*{Multiple roots}
Same-name channels must agree in promoted format, rate, epoch and subchannel
layout. Identical overlapping samples are accepted; inconsistent values are
rejected. The optimized streaming path falls back to the general overlap-aware
reader where required. The same rules apply to mixed HDF5 and raw roots.

\chapterpage{Writing an optional raw ring}
Raw mode is selected by constructing \api{raw::RawWriter}; ordinary readers and
recordings do not switch formats implicitly. This runnable example creates a
disposable 50 MS/s complex-i8 channel, publishes four samples and reads them
through the normal reader API. The example uses the development dependency
\api{tempfile = "3"}; applications can supply their own new channel directory.
\lstinputlisting{../examples/raw_ring.rs}
\textbf{In an acquisition loop}, pass each actual received block with its
absolute first sample index. Bytes are little-endian and interleaved in
(sample, subchannel, component) order. Complex components are real then
imaginary. For one complex-i8 channel, two bytes represent one sample. Forward
gaps are allowed and indexed; overlapping or backward writes are rejected.
Unaligned byte counts and sample-index overflows are errors.

\textbf{Publication and flushing.} Crossing a file cadence finalizes the old
chunk automatically. \api{flush()} explicitly finalizes a short current chunk;
you cannot append to that cadence interval afterwards. Dropping the writer
leaves the pending chunk unpublished. In a continuous loop, do not flush each
hardware packet; let cadence rollover publish complete chunks, then flush once
at a controlled stop.

\chapterpage{Raw options, format and operational limits}
\par\medskip
\begin{tabularx}{\linewidth}{@{}p{.43\linewidth}lX@{}}\toprule
\api{raw::Options} field & Default & Meaning\\\midrule
\api{component} & \api{I8} & \api{I8}, \api{I16}, \api{F32} or \api{F64}.\\
\api{complex} & \api{true} & Real/imaginary pairs; false for real samples.\\
\api{num_subchannels} & 1 & Positive number of interleaved subchannels.\\
\api{sample_rate_numerator} & 25,000,000 & Positive rational-rate numerator.\\
\api{sample_rate_denominator} & 1 & Positive rational-rate denominator.\\
\api{file_cadence_millisecs} & 1000 & Payload interval in milliseconds.\\
\api{subdir_cadence_secs} & 3600 & Time-directory interval in seconds.\\
\api{ring_files} & \api{None} & Retain archive; \api{Some(n)} keeps at most $n$ finalized chunks.\\
\api{durable} & \api{false} & Sync payload, metadata and commit directory before returning.\\\bottomrule
\end{tabularx}

Cadences must be positive; file cadence must divide the subdirectory cadence
and contain an integer number of samples. A ring size of zero is rejected.
\api{raw::Container::default()} is \api{Hdf5}; \api{Raw} is the explicit
alternative label. This enum is descriptive: it is not a writer factory.
\api{Component::bytes()} returns component size in bytes.

\subsection*{Container version 1}
The channel's \api{drf_properties.h5} contains ordinary rate/layout properties
and \api{rdigital_rf_raw_version = 1}. Each finalized chunk has a binary payload
\api{rf@SECONDS.MMM.raw} and an HDF5 sidecar \api{rf@SECONDS.MMM.h5} in the
usual time directory. Version 1 defines little-endian row-major bytes, without
sample scaling or compression. The sidecar stores:
\begin{itemize}
\item Attributes \api{rdigital_rf_raw_version}, \api{raw_component} (1=i8,
2=i16, 3=f32, 4=f64), \api{raw_is_complex}, \api{raw_num_subchannels} and \api{raw_rows}.
\item Dataset \api{rf_data_index}, an unsigned 64-bit $(n,2)$ array of absolute
sample starts and corresponding row offsets into the payload.
\end{itemize}
The payload is renamed into place first; the sidecar is published last as its
commit marker. Readers ignore partial files and verify payload length and
indices. No HDF5 sample dataset or per-sample HDF5 conversion is involved.

\subsection*{Ring lifetime and restart}
Chunks are immutable; pruning unpublishes the oldest metadata first, then
removes its payload. Open readers can retain an evicted inode until cache
removal; a new reader sees it as missing. The bound concerns finalized chunks,
not cached inodes or the pending chunk. At 50 MS/s, one complex-i8 channel
produces 100 MB/s; 40 one-second chunks are about 4 GB before overhead.

One writer is allowed per channel. Version 1 requires a new channel directory;
there is no append/resume constructor. After a crash, inspect the lock and
partial files before starting in a new directory. A RAM ring normally uses
\api{durable = false}; a timing benchmark of cached writes is not a durability
or sustained-disk benchmark. Ordinary HDF5 live-writer/SWMR consistency is not
guaranteed. The implementation currently excludes Digital Metadata, remote
transports and full Python API compatibility.

\smallskip
\textbf{Validation and further reference.}
Run \api{cargo test --all-targets}, \api{cargo test --doc} and
\api{cargo fmt --all -- --check}. The tests include gaps, boundaries, numeric
extremes, subchannels, mixed roots, partial chunks and ring eviction.
The upstream Python comparison covers a tested HDF5 subset, not the raw
extension. See \url{https://github.com/jvierine/rdigital_rf} and
\url{https://github.com/MITHaystack/digital_rf}.

\chapterpage{Benchmarks across reader and container options}
These measurements compare eight configurations with two access patterns. All
return the same complex-f32 values for a given stored component type, and every
output was checked exactly. The table gives median elapsed milliseconds over
five timed repeats, after one untimed warm-up. Lower is faster.

\par\medskip
% Numerical report and source provenance are retained in HDF5 metadata.
\begin{tabularx}{\linewidth}{@{}Xllrr@{}}\toprule
Container / cache & Type & API & Full & Sparse\\\midrule
HDF5, uncompressed & i8 & direct & 2.60 & 0.69 \\
HDF5, uncompressed & i8 & vector & 2.53 & 0.71 \\
HDF5, gzip-1 & i8 & direct & 52.25 & 3.86 \\
HDF5, gzip-1, 128 KiB cache & i8 & direct & 58.23 & 9.56 \\
HDF5, uncompressed & i16 & direct & 2.89 & 0.69 \\
Raw & i8 & direct & 2.03 & 0.13 \\
Raw & i8 & vector & 1.95 & 0.13 \\
Raw & i16 & direct & 2.71 & 0.14 \\
HDF5, uncompressed & i8 & i8 & 2.26 & 0.67 \\
Raw & i8 & i8 & 0.24 & 0.04 \\
\bottomrule
\end{tabularx}


\textbf{Access patterns.} Each fixture contains 2,097,152 complex samples in
one channel, one subchannel and one one-second interval. ``Full'' reads all
samples in one request. ``Sparse'' reads 32 windows of 4096 samples, spaced
65,536 samples apart: 131,072 samples in total. Compare configurations within
one column; the two columns transfer different amounts of data.

\textbf{API and cache settings.} ``Direct'' means \api{read_complex32};
``vector'' means \api{read_vector} with one subchannel and complex-f32 output,
without an additional conversion to complex-f64. Uncompressed HDF5 datasets
use contiguous storage. Gzip level 1 uses chunks of 65,536 complex-i8 samples.
The default HDF5 cache has a 64 MiB minimum; the capped case sets a 128 KiB
per-file ceiling. The same reader is retained across warm-up and timed repeats.

\textbf{Data and hardware.} The synthetic I/Q components are deterministic,
compressible staircase sequences in the signed 8-bit range. The i16 fixture
multiplies those values by 256. The combined payload and metadata occupy about
4.01 MiB for i8 HDF5 or raw, 8.01 MiB for i16, and 0.75 MiB for gzip-i8 HDF5.
The benchmark ran on an Intel Core i9-11950H Linux x86-64 system with ext4 and a
warmed OS page cache. Timing excludes fixture construction and correctness
checks, but includes allocation of the returned vectors.

\textbf{Interpretation.} Raw is about 2.8 times faster than the default gzip
case for full reads here. Warm, contiguous, uncompressed HDF5 is about 6.3 times
faster than raw for the same i8 selection. The direct API is about 8.9 times
faster than the general vector API for that uncompressed HDF5 full read. A
larger HDF5 cache improves repeated sparse compressed reads by about 2.7 times
in this workload. Raw therefore does not provide a universal speed advantage;
choose the container, API and cache policy using the intended access pattern.

These are warm-cache read measurements, not sustained disk, acquisition,
writer durability or end-to-end receiver throughput. Compressibility is
synthetic and does not predict the size of radio-noise recordings. A separate
40-million-sample read on DOB took 1.10 s from its HDF5 recording and 0.42 s
from raw. That different recording, hardware and workload is not directly
comparable with this table. Numerical reports retain every repeat, extrema,
fixture size and correctness result; reproduce the matrix on the target system.
% Reproduce with cargo run --release --example bench_features -- NEW_REPORT.h5.
% Benchmark source and toolchain provenance are in the HDF5 report metadata.
\end{document}
